Investor comments on financial social media contain rich sentiment information, and identifying their sentiment accurately is fundamental to financial public-opinion monitoring and investor-sentiment research. Large language models (LLMs) excel at text understanding, and parameter-efficient fine-tuning methods such as LoRA make it possible to fine-tune them on a single consumer GPU. However, whether their advantage over traditional methods in Chinese financial sentiment classification is reliable, and whether their confidence scores can be trusted, requires evaluation under explicit data-quality and statistical assumptions.

Using FinFE, a three-class Chinese financial social-media sentiment dataset, this thesis compares TF-IDF with logistic regression and support vector machines, fine-tuned BERT and RoBERTa-wwm-ext, and zero-shot prompting, few-shot prompting and LoRA fine-tuning of three open-source LLMs (Qwen3-8B, Qwen2.5-7B-Instruct and GLM-4-9B) on a single RTX 3090 GPU. The results are analysed from three perspectives: statistical inference, confidence reliability and data quality.

The main findings are as follows. (1) About 54\% of the original evaluation records share normalized text with the original training file. Training exposure is distinguished from validation exposure; cleaning removes exact cross-split duplicates, and cleaning strategies are compared on shared unexposed examples. (2) LoRA-tuned LLMs achieve a macro-F1 of 0.798--0.807, compared with 0.778 for RoBERTa and about 0.70 for TF-IDF. Zero-shot and few-shot prompting perform worse than supervised fine-tuning. (3) The LLMs have higher F1 point estimates across the three training seeds on the fixed data split, with significant post-hoc comparisons in the multi-seed analysis. However, for seed 42, McNemar tests and paired bootstrap intervals do not establish a significant advantage over RoBERTa after multiplicity adjustment. (4) Zero-shot Qwen3 has a mean confidence of 96.4\% but an accuracy of 67.1\%. Five-fold temperature cross-fitting within the test distribution gives LoRA models an expected calibration error of 0.013--0.030; accuracy is about 96\% when answering only the most confident half of the examples. Standard conformal prediction achieves about 90\% overall set coverage, with 77\%--78\% singleton sets. Accuracy within the singleton subset is 88.7\%--89.3\%, not a separately guaranteed 90\%. (5) Duplicate-record pairs agree on 82.9\% of labels, and 82\% of conflicts involve the neutral class. Without evidence of independent reannotation, a noise-model calculation is not interpreted as an empirical accuracy ceiling. For seed 42, 250 of Qwen3 LoRA's 307 errors involve the neutral class; deterministically sampled cases illustrate boundaries without causal attribution. In one run, 4-bit QLoRA performs comparably to 16-bit LoRA, with 11.2 GB of peak framework-allocated GPU memory.

These results suggest that evaluations of LLMs should jointly consider data quality, statistical significance and confidence reliability, and they provide a reference for model selection and human--machine collaboration in financial public-opinion monitoring.

\textbf{Keywords:} large language models; LoRA fine-tuning; financial sentiment analysis; statistical testing; confidence calibration
